\documentclass[10pt]{article}
\usepackage[margin=1in]{geometry}
\usepackage[T1]{fontenc}
\usepackage{lmodern,microtype,amsmath,amssymb,amsthm}
\theoremstyle{definition}
\newtheorem{assumption}{Assumption}
\begin{document}
\section*{Quasar Optimization with Amortized Linear Oracles}
We consider minimizing a differentiable function $f:E\to\mathbb{R}$ over a convex set $C$ in a real Hilbert space $E$. Its gradient is $L$-Lipschitz, where $L\geq0$, and a minimizer $x^\star\in C$ satisfies $\tau(f(x)-f(x^\star))\leq\langle\nabla f(x),x-x^\star\rangle$ for all $x\in C$, with $\tau\in(0,1]$. The oracle returns exact gradients or approximate linear minimizers; linear replies at odd call indices have nonnegative additive errors totaling at most $\Gamma\geq0$. Given $\varepsilon>0$, the required output $\bar{x}\in C$ satisfies $f(\bar{x})-f(x^\star)\leq\varepsilon$ almost surely. Only $E$, $L$, $\tau$, and an initial point $x^0\in C$ are public inputs.

The constrained optimization problem is
\[
\min_{x\in C} f(x),
\]
where the norm on $E$ is induced by its inner product. The feasible set, objective, and oracle satisfy the following conditions.

\begin{assumption}[Feasible geometry]
For a parameter $D\geq0$, the set $C\subseteq E$ is convex, contains $x^0$, and satisfies
\[
\|x-y\|\leq D
\qquad\text{for every }x,y\in C.
\]
\end{assumption}

\begin{assumption}[Global smoothness]
The function $f$ is differentiable everywhere on $E$, with gradient $\nabla f(x)$ at each $x\in E$, and
\[
\|\nabla f(x)-\nabla f(y)\|\leq L\|x-y\|
\qquad\text{for every }x,y\in E.
\]
\end{assumption}

\begin{assumption}[Attainment and quasar inequality]
There exists $x^\star\in C$ such that, for every $x\in C$,
\[
\begin{aligned}
f(x^\star)&\leq f(x),\\
\tau\bigl(f(x)-f(x^\star)\bigr)
&\leq\langle\nabla f(x),x-x^\star\rangle.
\end{aligned}
\]
\end{assumption}

\begin{assumption}[Exact gradients and amortized linear minimization]
Oracle calls have indices $j=0,1,2,\ldots$. Each call specifies either a gradient query at $x\in E$ or a linear-minimization query with direction $v\in E$, and returns one vector in $E$. At every call index, a gradient query returns $\nabla f(x)$ for every realization of the oracle randomness.

There exists a deterministic sequence $(\delta_k)_{k\geq0}$ satisfying
\[
\begin{aligned}
\delta_k&\geq0 &&\text{for every }k\geq0,\\
\sum_{k=0}^{n-1}\delta_k&\leq\Gamma
&&\text{for every }n\geq0.
\end{aligned}
\]
At call index $j=2k+1$, the reply $s\in E$ to any linear-minimization query $v\in E$ satisfies, for every realization of the oracle randomness,
\[
\begin{aligned}
s&\in C,\\
\langle v,s\rangle&\leq\langle v,u\rangle+\delta_k
\qquad\text{for every }u\in C.
\end{aligned}
\]
These linear-minimization guarantees apply only at odd indices in the total sequence of calls. Each query of either kind costs one oracle call. Any internal algorithm randomness is independent of the oracle randomness.
\end{assumption}

Queries may be chosen adaptively, with gradient query points and linear query directions anywhere in $E$. The algorithm is selected using only $(E,L,x^0,\tau)$; the set $C$ and parameters $D$, $\Gamma$, and $\varepsilon$ are not supplied to it. An externally specified budget $K\in\mathbb{N}$ may depend on the full parameter tuple $(E,L,x^0,\tau,C,D,\Gamma,\varepsilon)$ and counts both gradient and linear-minimization calls.

After $K$ calls, the algorithm returns $\bar{x}\in E$. Writing $\mathbb{P}$ for the joint law of the algorithm and oracle randomness, success requires a minimizer $x^\star\in C$ for which
\[
\bar{x}\in C
\quad\text{and}\quad
f(\bar{x})-f(x^\star)\leq\varepsilon
\qquad\mathbb{P}\text{-almost surely}.
\]
The minimizer in the quasar inequality may be used here, since all minimizers over $C$ have the same objective value.
\end{document}
